% Migrado de legacy/Hidrología/Examenes/Parcial 1/ (parcial_1_hidrologia.tex
% + parcial_1_hidrologia_questions.tex) -- examen real completo de
% selección múltiple, 32 reactivos en 3 unidades, sin recortar.
%
% SIN \DocumentMetadata a propósito -- ver STYLE.md § PDF etiquetado
% (el entorno \begin{solution} de exam.cls choca con tagpdf). Este
% ejemplo no usa \begin{solution} porque el original resuelve las
% respuestas correctas con \CorrectChoice (visibles solo con la opción
% de clase "answers"), no con explicaciones de solución aparte.
%
% Para generar la clave de respuestas del mismo archivo:
%   \documentclass[11pt,addpoints,answers]{exam}
\documentclass[11pt,addpoints]{exam}
\usepackage{icv}

% -----------------------------------------------------------------------
% Metadata
% -----------------------------------------------------------------------
\icvsetup{
  asignatura  = {ICV 342-T -- Hidrología},
  titulocorto = {Examen Parcial 1 -- Hidrología},
  titulo      = {Examen Parcial 1},
  subtitulo   = {Unidades I--III: introducción, meteorología, cuencas hidrográficas},
  profesor    = {Dr. Ricardo Hernández Moreira},
  periodo     = \icvciclo{2026}{3},
  version     = {v1},
}

% -----------------------------------------------------------------------
% Localización de exam.cls al español
% -----------------------------------------------------------------------
\pointpoints{punto}{puntos}
\vpword{Puntos}
\vsword{Calificación}
\vtword{Total:}
\hpword{Puntos:}
\hsword{Calificación:}
\htword{Total}
\qformat{\textbf{Pregunta \thequestion.} \hfill (\thepoints)}
\pointsinrightmargin
\setlength{\rightpointsmargin}{1.5cm}

\begin{document}
\icvmaketitle

\begin{center}
\begin{tabular}{ll}
  \toprule
  \textbf{Fecha:} & Fecha del examen \\
  \textbf{Duración:} & 100 minutos \\
  \textbf{Puntaje total:} & \numpoints\ puntos \\
  \bottomrule
\end{tabular}
\end{center}

\vspace{6pt}
\begin{center}
\makebox[\linewidth][l]{%
  \textbf{Nombre:}\hspace{0.5em}\rule{0.55\linewidth}{0.4pt}%
  \hspace{1em}\textbf{ID:}\hspace{0.5em}\hrulefill%
}
\end{center}
\vspace{6pt}

\begin{icvnote}[Instrucciones]
Selección múltiple. Marque la mejor respuesta para cada reactivo. Examen
cerrado -- no se permite material de consulta ni calculadora, salvo
donde se indique lo contrario.
\end{icvnote}

\begin{questions}

\item[]\textbf{Unidad I. Introducción, ciclo hidrológico y contexto nacional}\par\medskip

\question[2] La hidrología se define mejor como el estudio de:
\begin{choices}
  \choice la composición química de los océanos y mares.
  \choice la distribución espacial de la atmósfera terrestre.
  \CorrectChoice el agua en la Tierra, su ocurrencia, circulación y propiedades.
  \choice la formación exclusiva de nubes y tormentas.
\end{choices}

\question[2] La idea central del ciclo hidrológico es que:
\begin{choices}
  \choice el agua se produce continuamente por evaporación desde los océanos.
  \choice el agua permanece fija en un solo compartimiento natural.
  \CorrectChoice el agua circula entre atmósfera, superficie, subsuelo y océanos con intercambios continuos.
  \choice la precipitación ocurre solo en áreas montañosas.
\end{choices}

\question[3] Cuando se dice que un país puede tener agua en promedio anual
pero seguir sufriendo escasez, se está enfatizando principalmente:
\begin{choices}
  \choice la falta total de precipitación.
  \CorrectChoice la distribución espacial y temporal desigual del recurso.
  \choice que los ríos no tienen caudal base.
  \choice que la evaporación siempre es cero.
\end{choices}

\question[3] ¿Cuál de los siguientes problemas es típicamente hidrológico
y no solo administrativo?
\begin{choices}
  \choice Falta de etiquetas en los planos de una obra.
  \CorrectChoice Variabilidad extrema entre época húmeda y época seca.
  \choice Error tipográfico en un informe.
  \choice Cambio de formato en una presentación.
\end{choices}

\question[3] En la República Dominicana, una lectura básica del problema
del agua debe incluir:
\begin{choices}
  \choice solo el uso doméstico urbano.
  \choice solo el control de inundaciones.
  \CorrectChoice disponibilidad, demandas, calidad, variabilidad climática y gestión institucional.
  \choice únicamente la desalinización.
\end{choices}

\question[3] El Plan Hidrológico Nacional 2025--2045 presenta como idea
general:
\begin{choices}
  \choice una colección de datos meteorológicos sin interpretación de gestión.
  \CorrectChoice una estrategia de planificación hídrica para seguridad hídrica y gestión integrada.
  \choice un documento solo para diseño de presas.
  \choice un inventario exclusivo de pozos.
\end{choices}

\question[2] ¿Qué capítulo del Plan Hidrológico Nacional trata el marco
físico y climatológico?
\begin{choices}
  \choice Capítulo 1.
  \CorrectChoice Capítulo 2.
  \choice Capítulo 5.
  \choice Capítulo 6.
\end{choices}

\question[2] ¿Qué capítulo del Plan Hidrológico Nacional trata la
caracterización del ámbito del estudio?
\begin{choices}
  \choice Capítulo 2.
  \choice Capítulo 4.
  \CorrectChoice Capítulo 5.
  \choice Capítulo 7.
\end{choices}

\question[2] ¿Qué capítulo del Plan Hidrológico Nacional desarrolla
recursos hídricos, demandas y balances?
\begin{choices}
  \choice Capítulo 1.
  \choice Capítulo 3.
  \choice Capítulo 5.
  \CorrectChoice Capítulo 6.
\end{choices}

\question[3] Según la lectura general del PHN, una prioridad operativa
coherente con el documento es:
\begin{choices}
  \choice aumentar solo el consumo sin controlar pérdidas.
  \CorrectChoice mejorar conocimiento, disponibilidad, demanda, calidad y adaptación climática.
  \choice ignorar las cuencas transfronterizas.
  \choice sustituir toda gestión por análisis de laboratorio.
\end{choices}

\question[3] El ciclo hidrológico conecta de forma más directa:
\begin{choices}
  \choice solo clima y vegetación.
  \CorrectChoice precipitación, escurrimiento, infiltración, evaporación y almacenamiento.
  \choice solo aguas subterráneas y océanos.
  \choice solo radiación solar y nubosidad.
\end{choices}

\item[]\textbf{Unidad II. Meteorología e hidrología}\par\medskip

\question[2] La diferencia principal entre tiempo atmosférico y clima es
que:
\begin{choices}
  \choice el clima cambia por hora y el tiempo por décadas.
  \CorrectChoice el tiempo describe condiciones de corto plazo y el clima patrones de largo plazo.
  \choice ambos términos significan exactamente lo mismo.
  \choice el clima solo aplica al mar.
\end{choices}

\question[2] ¿Cuál variable meteorológica está más directamente asociada
con la energía disponible para evaporación?
\begin{choices}
  \choice Presión de vapor.
  \CorrectChoice Radiación solar.
  \choice Color del suelo.
  \choice Distancia al río.
\end{choices}

\question[2] En una estación meteorológica, el instrumento asociado más
directamente con la velocidad del viento es:
\begin{choices}
  \choice pluviómetro.
  \choice barómetro.
  \CorrectChoice anemómetro.
  \choice evaporímetro de bandeja únicamente.
\end{choices}

\question[2] La humedad relativa se interpreta como:
\begin{choices}
  \choice la masa total de lluvia en un mes.
  \CorrectChoice la razón entre vapor de agua presente y vapor máximo posible a esa temperatura.
  \choice la cantidad de nieve acumulada.
  \choice la diferencia entre presión y temperatura.
\end{choices}

\question[2] Un aumento de temperatura, manteniendo lo demás constante,
tiende a:
\begin{choices}
  \choice reducir siempre la evaporación.
  \CorrectChoice favorecer la capacidad de evaporación del aire.
  \choice eliminar la precipitación.
  \choice impedir la transpiración.
\end{choices}

\question[3] El vínculo correcto entre meteorología e hidrología es que:
\begin{choices}
  \choice los procesos meteorológicos no afectan el agua superficial.
  \CorrectChoice el clima y el tiempo controlan entradas, pérdidas y variabilidad del ciclo hidrológico.
  \choice solo la presión atmosférica explica los caudales.
  \choice la hidrología no usa datos meteorológicos.
\end{choices}

\question[3] Si se desea estudiar el impacto de un cambio climático sobre
recursos hídricos, una variable básica a revisar es:
\begin{choices}
  \choice la marca del pluviómetro.
  \CorrectChoice la tendencia de precipitación y temperatura.
  \choice el color de los gráficos.
  \choice el nombre del laboratorio.
\end{choices}

\question[3] La evaporación se considera en hidrología como:
\begin{choices}
  \choice una ganancia de agua para la cuenca.
  \CorrectChoice una pérdida de agua hacia la atmósfera.
  \choice un tipo de escurrimiento subterráneo.
  \choice una medida exclusiva de las presas.
\end{choices}

\question[3] Una relación útil entre meteorología e hidrología en el
curso es:
\begin{choices}
  \choice la precipitación no se relaciona con escurrimiento.
  \CorrectChoice la precipitación alimenta la cuenca y se reparte entre almacenamiento, escurrimiento e infiltración.
  \choice la evaporación solo ocurre en invierno.
  \choice el viento no interviene en procesos hídricos.
\end{choices}

\item[]\textbf{Unidad III. Cuencas hidrográficas}\par\medskip

\question[3] Una cuenca hidrográfica es:
\begin{choices}
  \choice cualquier área urbana.
  \CorrectChoice el área que drena escorrentía hacia un punto de salida común.
  \choice una zona sin drenaje superficial.
  \choice un acuífero confinado.
\end{choices}

\question[3] El divisor topográfico de aguas se refiere a:
\begin{choices}
  \choice la línea de máxima contaminación.
  \CorrectChoice la línea que separa cuencas vecinas.
  \choice el cauce principal.
  \choice el nivel de evaporación.
\end{choices}

\question[3] La importancia de la cuenca como unidad de análisis
hidrológico radica en que:
\begin{choices}
  \choice simplifica la meteorología global.
  \CorrectChoice integra procesos de entrada, almacenamiento, transferencia y salida del agua.
  \choice elimina la necesidad de medir precipitación.
  \choice solo sirve para inventarios forestales.
\end{choices}

\question[3] Entre los rasgos fisiográficos que influyen en la respuesta
de la cuenca está:
\begin{choices}
  \CorrectChoice la forma, el relieve y la red de drenaje.
  \choice solo el color del suelo.
  \choice la edad del ingeniero.
  \choice el formato del informe.
\end{choices}

\question[3] En términos morfométricos, la respuesta hidrológica de una
cuenca está controlada principalmente por:
\begin{choices}
  \choice la forma y el relieve.
  \choice la red de drenaje.
  \CorrectChoice A y B son correctas.
  \choice solo el color del suelo.
\end{choices}

\question[3] La densidad de drenaje se relaciona mejor con:
\begin{choices}
  \choice la suma del nombre de los ríos.
  \CorrectChoice la longitud de cauces por unidad de área.
  \choice la evaporación diaria.
  \choice la profundidad del acuífero.
\end{choices}

\question[3] Si una cuenca tiene \qty{36}{\kilo\meter} de cauces y
\qty{12}{\kilo\meter\squared} de área, su densidad de drenaje es:
\begin{choices}
  \choice \qty{0.33}{\kilo\meter\per\kilo\meter\squared}.
  \choice \qty{2}{\kilo\meter\per\kilo\meter\squared}.
  \CorrectChoice \qty{3}{\kilo\meter\per\kilo\meter\squared}.
  \choice \qty{4}{\kilo\meter\per\kilo\meter\squared}.
\end{choices}

\question[3] Una cuenca alargada, en comparación con una compacta y con
lluvia similar, tiende a producir:
\begin{choices}
  \choice siempre el mismo pico.
  \CorrectChoice un pico más atenuado y menos concentrado en el tiempo.
  \choice cero escurrimiento.
  \choice mayor evaporación instantánea siempre.
\end{choices}

\question[3] Una cuenca con área de \qty{50}{\kilo\meter\squared} y
longitud principal de \qty{10}{\kilo\meter} tiene un factor de forma
aproximado de:
\begin{choices}
  \choice 0.05.
  \choice 0.2.
  \CorrectChoice 0.5.
  \choice 2.0.
\end{choices}

\question[3] ¿Cuál opción diferencia mejor escurrimiento superficial y
subsuperficial?
\begin{choices}
  \choice ambos son exactamente iguales.
  \CorrectChoice el superficial circula sobre el terreno y el subsuperficial circula dentro de capas del suelo someras.
  \choice el subsuperficial solo ocurre en el océano.
  \choice el superficial depende únicamente de pozos.
\end{choices}

\question[3] Un uso frecuente de la relación $Q = A \cdot v$ es estimar:
\begin{choices}
  \choice precipitación media areal.
  \CorrectChoice caudal a partir de área hidráulica y velocidad media.
  \choice humedad relativa.
  \choice almacenamiento atmosférico.
\end{choices}

\question[3] ¿Qué efecto suele tener la urbanización sobre la respuesta
hidrológica de una cuenca?
\begin{choices}
  \choice reduce siempre el gasto pico.
  \CorrectChoice aumenta la rapidez de respuesta y tiende a elevar el pico de escurrimiento.
  \choice elimina los cauces.
  \choice convierte la lluvia en evaporación.
\end{choices}

\question[3] Una función ambiental importante del bosque en la cuenca es:
\begin{choices}
  \choice aumentar la erosión.
  \CorrectChoice ayudar a proteger suelos, regular escorrentía y mejorar la infiltración.
  \choice impedir toda precipitación.
  \choice secar los acuíferos de forma inmediata.
\end{choices}

\question[3] La principal diferencia entre cuenca superficial y cuenca
subterránea es que:
\begin{choices}
  \choice la primera existe solo en planos y la segunda solo en mapas.
  \CorrectChoice la primera se define por divisores topográficos y la segunda por divisores piezométricos o hidrogeológicos.
  \choice ambas son idénticas.
  \choice la segunda solo aparece en desiertos.
\end{choices}

\question[8] Explique en 4 a 6 líneas cómo la forma de la cuenca y su red
de drenaje influyen en el tiempo al pico y en el gasto pico del
hidrograma.

\vspace{3cm}

\end{questions}

\end{document}
